% !TeX TS-program = xelatex
% 虚构演示简历 — 仅用于 hellojobs 测试数据
\documentclass{resume-photo}
\usepackage{xeCJK}
\setCJKmainfont{Noto Serif CJK SC}
\usepackage{fontspec}
\usepackage{enumitem}
\setlist[itemize]{itemsep=0pt, topsep=1pt, parsep=0pt, partopsep=0pt}

\ResumeName{赵磊}

\begin{document}

\ResumeContacts{
  138-0006-0006,%
  \ResumeUrl{mailto:zhaolei@hust.edu.cn}{zhaolei@hust.edu.cn},%
  \textnormal{华中科技大学 | 软件工程 · 学士 | 2020—2024}%
}

\ResumeTitle

\noindent Android 客户端开发，熟悉 Kotlin、Jetpack Compose 与性能剖析；主导包体积治理，APK 减小 18\%。注重可观测性与文档沉淀，习惯在评审阶段拆解风险；参与线上复盘与跨团队联调，英语 CET-6，可阅读官方文档与 RFC（演示数据）。

\section{教育背景}
\ResumeItem
{华中科技大学}
[\textnormal{软件学院 | 软件工程 | 学士}]
[2020.09—2024.06]

\begin{itemize}
  \item 毕业设计：短视频客户端弱网优化。
\end{itemize}


\ResumeItem
{实验室与助教}
[\textnormal{本科生科研 / 课程助教}]
[2022—2025]

\begin{itemize}
  \item 进入校级重点实验室参与课题，负责数据采集、清洗与可视化看板搭建。
  \item 担任《数据结构》《操作系统》课程助教，批改作业 200+ 份，答疑课时 40h+。
  \item 组织期中复习串讲与上机辅导，课程评教得分位于前 10\%（演示）。
  \item 整理课程 FAQ 与实验踩坑文档，被后续年级沿用。
\end{itemize}



\section{专业技能}
\begin{itemize}
  \item 熟悉 Kotlin、Android SDK、MVVM 与协程并发。
  \item 掌握 LeakCanary、Systrace 与电量/卡顿治理。
  \item 了解 JNI 与音视频硬解基础。
  \item 熟悉 Linux 日常运维与 Shell，具备日志检索与 CPU/内存突增初步定位能力。
  \item 掌握 Git / CI 与 Code Review，了解 Docker 与 Kubernetes 基本编排与滚动发布。
  \item 了解 OAuth2 / JWT、接口幂等与 REST 规范，具备单元测试与集成测试习惯（JUnit/pytest 等）。
  \item 能阅读英文技术文档，具备跨团队沟通与需求澄清能力（测试简历）。
\end{itemize}

\section{实习经历}
\ResumeItem
{小红书}
[\textnormal{Android 开发实习生}]
[2023.12—2024.05]

\begin{itemize}
  \item \textbf{职责}: 笔记发布链路稳定性与启动耗时治理。
  \item \textbf{成果}: 冷启动主线程耗时降低 120ms（P90）。
\end{itemize}


\ResumeItem
{贝壳找房 | 如视}
[\textnormal{C++/Go 工程实习生}]
[2023.07—2024.01]

\begin{itemize}
  \item \textbf{职责}: 参与三维空间数据处理流水线，负责任务队列与重试策略。
  \item \textbf{优化}: 调整 gRPC 流式传输参数，大文件传输失败重试率下降 45\%。
  \item \textbf{质量}: 引入 fuzz 测试覆盖解析模块，发现边界崩溃 3 类并修复。
  \item \textbf{部署}: 熟悉 Ansible 批量发布与蓝绿切换流程。
  \item \textbf{成长}: 完成从业务代码到性能剖析工具链的完整闭环实践。
\end{itemize}



\section{项目经历}
\ResumeItem
{Compose 实验性 UI 套件}
[\textnormal{个人开源}]
[2023.06—2023.12]

\begin{itemize}
  \item GitHub 300+ stars，被 2 个示例项目引用。
\end{itemize}


\ResumeItem
{多租户 SaaS 计费引擎}
[\textnormal{订阅制 + 用量计费（演示）}]
[2023.09—2024.02]

\begin{itemize}
  \item \textbf{模型}: 设计套餐、用量累计、账单周期与 proration 规则。
  \item \textbf{一致}: 使用 Outbox + 消息表保证计费事件与财务对账一致。
  \item \textbf{性能}: 批价任务分片并行，账单生成时间缩短 48\%。
  \item \textbf{测试}: 构造边界用例库 150+，覆盖闰月与升级降级场景。
\end{itemize}



\section{个人荣誉}
\begin{itemize}
  \item 蓝桥杯软件类 全国一等奖
  \item 华中科技大学 自强奖学金
  \item 全国计算机等级考试四级（网络工程师）（演示）
  \item 校级「优秀学生干部 / 优秀团员」或技术社团骨干（综合测评前 15\%）
  \item 开源贡献：向知名项目提交被合并的 PR（文档/测试/feature）（演示）
\end{itemize}

\end{document}
